\figura[0.45]{sol_fig1.png}

\begin{solapartat}
\punts{0,2} Mostra la variació de la freqüència amb la velocitat de la font.\\
$\rightarrow$ \underline{\textbf{Efecte Doppler}}

\punts{0,2} Segons la gràfica, una font amb $v = +100\ \mathrm{m/s}$ la sentim amb $f = 240\ \mathrm{Hz}$ i una font amb $v = -100\ \mathrm{m/s}$ amb $f = 130\ \mathrm{Hz}$. La freqüència més alta ha de correspondre a la font acostant-se a nosaltres, per tant el canvi serà:

\punts{0,4} $f_\mathit{final} - f_\mathit{inicial} = 130 - 240 = -110\ \mathrm{Hz}$,

\punts{0,2} La freqüència \underline{\textbf{disminuirà.}}
\end{solapartat}

\begin{solapartat}
\punts{0,1} $L_I = 10\log\dfrac{I}{I_0}$

\punts{0,2} $I = I_0\,10^{\frac{L_I}{10}} = (1\times 10^{-12})\times 10^{6,5} = 3,2\times 10^{-6}\ \mathrm{W/m^2}$ \quad \punts{0,1}

\punts{0,2} Aquesta intensitat correspon a 25~m, tal com es veu a la taula.

\punts{0,2} $\textit{Potència} = I\times 4\pi R^2$

\punts{0,2} $P = 3,2\times 10^{-6}\times 4\pi(25)^2 = 0,025\ \mathrm{W}$
\end{solapartat}
